\section{27.02.2026}{Sploty}

Przyjmujemy oznaczenia:
\begin{itemize}
  \item $G$ - grupa topologiczna
  \item $\Hh$ - przestrzeń Hilberta
  \item $\U(\Hh)$ - unitarne przekształcenia $\Hh$.
\end{itemize}

\begin{definition}{topologia na $\B(\Hh)$}{słaba/silna topologia}
  Niech $\B(\Hh)$ będzie przestrzenią operatorów ograniczonych na $\Hh$. Definiujemy wtedy dwie topologie na $\B(\Hh)$:
  \begin{description}
    \item[silną] (piszemy SO) jako najmniejszą topologię, w której dla każdego $v\in\Hh$ odwzorowania $g\mapsto gv$ są ciągłe,
    \item[słabą] (piszemy WO) jako najmniejszą topologię, w której dla wszystkich $v,w\in\Hh$ odwzorowania $g\mapsto \langle v,gw\rangle\in\C$ są ciągłe.
  \end{description}
\end{definition}

WO zachowuje się dobrze na liniowych przestrzeniach topologicznych, tzn. jeśli $V$ jest przestrzenią liniową, to dla wszystkich $v\in V$ oraz $\xi\in V^*$ odwzorowanie $g\mapsto \xi(gv)$ jest ciągłe.

\begin{remark}{}{}
  Odwzorowanie $G\to \U(\Hh)$ jest ciągle względem SO/WO $\iff$ $(\forall\;g\in G)\;g\mapsto gv/g\mapsto \langle w,gv\rangle$ są ciągłe.
\end{remark}

\begin{fact}{}{}
  SO i WO topologie na $\U(\Hh)$ są takie same, ale nie jest to prawdą na $\B(\Hh)$
\end{fact}

\begin{proof}
  Zadanie.
\end{proof}





